\documentclass[12pt]{article}
\usepackage{amsmath,amssymb}
\usepackage{geometry}
\usepackage{enumitem}
\geometry{margin=0.8in}
\setlength{\parindent}{0pt}

\begin{document}

\begin{center}
  {\Large\textbf{Trigonometric Equation Solving}}\\[3pt]
  {\large Algebra II / Precalculus Practice}
\end{center}

\noindent Name: \rule{2.7in}{0.4pt}\hfill Date: \rule{1.5in}{0.4pt}

\section*{Instructions}
Complete all 12 problems in 45 minutes without a calculator. Show the algebra that produces each trigonometric equation and give exact answers. In Problems 1--6, solve on $0\le x<2\pi$. In Problems 7--12, give all real solutions using $k\in\mathbb Z$. If an angle is not a standard unit-circle value, leave it in inverse-trigonometric form.

Use only factoring, inverse trigonometric functions, the complementary-angle rules, and the identities
\[
\sin^2x+\cos^2x=1,
\qquad
1+\tan^2x=\sec^2x,
\qquad
1+\cot^2x=\csc^2x.
\]
Do not use sum, difference, double-angle, half-angle, or product-to-sum identities. Expressions such as $\cos(2x)$ should be treated as cosine with argument $2x$.

For equations with matching trigonometric functions, the following general patterns may be used:
\[
\begin{aligned}
\sin A=\sin B
&\Longrightarrow A=B+2\pi k
\quad\text{or}\quad
A=\pi-B+2\pi k,\\
\cos A=\cos B
&\Longrightarrow A=2\pi k\pm B,\\
\tan A=\tan B
&\Longrightarrow A=B+\pi k.
\end{aligned}
\]

\section*{Solve on $0\le x<2\pi$}
\begin{enumerate}[leftmargin=*,itemsep=1.35em]
  \item
  \[
    2\sin x+\sqrt2=0
  \]

  \item
  \[
    2\cos^2x-3\cos x+1=0
  \]

  \item
  \[
    \tan^2x=3
  \]

  \item
  \[
    2\sin^2x+\cos x-2=0
  \]

  \item
  \[
    \sec^2x-3\tan x-1=0
  \]

  \item
  \[
    \cos(2x)=\sin x
  \]
\end{enumerate}

\section*{Give all real solutions}
\begin{enumerate}[leftmargin=*,itemsep=1.35em,resume]
  \item
  \[
    \sin x=\frac25
  \]

  \item
  \[
    \cos(2x)=\cos(5x)
  \]

  \item
  \[
    \sin(3x)=\cos x
  \]

  \item
  \[
    3\csc^2x-4\cot^2x=2
  \]

  \item
  \[
    \tan(2x)=-\sqrt3
  \]

  \item
  \[
    \sin^{-1}(2\cos x)=-\frac{\pi}{6}
  \]
\end{enumerate}

\end{document}
